\chapter{Physiology}

\medterm{Active Transport}-- Movement of substances across cell membranes against their concentration gradient, requiring energy input, typically from ATP hydrolysis. Primary active transport directly uses ATP: the Na+/K+ ATPase pumps 3 Na+ out and 2 K+ in per ATP, maintaining the electrochemical gradient essential for neuronal signaling, cell volume regulation, and secondary active transport. Other examples include the Ca2+ ATPase, H+/K+ ATPase (gastric acid secretion), and ABC transporters. Secondary active transport uses the Na+ gradient established by the Na+/K+ ATPase to cotransport or countertransport other solutes (e.g., SGLT1 for glucose and amino acid cotransport in the intestine, Na+/Ca2+ exchanger in cardiac muscle).

\medterm{Aldosterone}-- A mineralocorticoid hormone produced by the zona glomerulosa of the adrenal cortex, the final effector of the renin-angiotensin-aldosterone system (RAAS). Aldosterone acts on principal cells of the collecting duct, increasing ENaC (epithelial sodium channel) expression on the apical membrane and Na+/K+ ATPase on the basolateral membrane, promoting sodium reabsorption and potassium secretion. Water follows sodium osmotically, expanding extracellular fluid volume and increasing blood pressure. Aldosterone secretion is stimulated by angiotensin II, hyperkalaemia, and ACTH, and inhibited by ANP and hypokalaemia. Excessive aldosterone (Conn syndrome) causes hypertension with hypokalaemia, metabolic alkalosis, and suppressed renin.

\medterm{Allosteric Regulation}-- The regulation of enzyme activity by the binding of an effector molecule at a site other than the active site (allosteric site), causing a conformational change that alters the enzyme's activity. Allosteric enzymes typically have multiple subunits with cooperative binding kinetics (sigmoidal rather than hyperbolic velocity curves). Positive allosteric effectors (activators) shift the curve leftward (increase affinity for substrate, decrease Km), while negative allosteric effectors (inhibitors) shift the curve rightward (decrease affinity, increase Km). Classic examples include oxygen binding to haemoglobin (cooperative, with 2,3-BPG as a negative allosteric effector), aspartate transcarbamoylase (ATCase) in pyrimidine synthesis, and feedback inhibition in metabolic pathways.

\medterm{Amino Acid Catabolism}-- The breakdown of amino acids after protein digestion, with the amino group removed and the carbon skeleton metabolized for energy or gluconeogenesis/ketogenesis. Transamination: amino groups are transferred to alpha-ketoglutarate, forming glutamate, by aminotransferases requiring pyridoxal phosphate (PLP, vitamin B6) as a cofactor. ALT (alanine aminotransferase) and AST (aspartate aminotransferase) are clinically important: elevated levels indicate liver damage (ALT is more liver-specific) or myocardial injury (AST). Following transamination, the alpha-keto acids enter gluconeogenesis or the TCA cycle, and ammonium is converted to urea in the urea cycle.

\medterm{Baroreceptor Reflex}-- The rapid neural mechanism for short-term blood pressure regulation, involving baroreceptors (mechanoreceptors) in the carotid sinus (innervated by CN IX) and aortic arch (innervated by CN X). When blood pressure rises, baroreceptors increase their firing rate, signaling the cardiovascular center in the medulla. The response: increased parasympathetic (vagal) output to the heart (decreasing HR) and decreased sympathetic output (decreasing HR, contractility, and peripheral resistance), lowering blood pressure towards normal. The reflex is essential for maintaining blood pressure stability during postural changes and blood volume fluctuations.

\medterm{Basal Metabolic Rate (BMR)}-- The minimum energy expenditure required to maintain vital physiological functions at rest in a thermoneutral environment, measured after 12 hours of fasting and 8 hours of sleep. BMR accounts for approximately 60--75% of total daily energy expenditure in sedentary individuals. It is influenced by body composition (lean mass is the primary determinant), age (declines ~1--2% per decade after age 20), sex (5--10% higher in males), thyroid status, and genetics. BMR is estimated using predictive equations (Harris-Benedict, Mifflin-St Jeor) or measured by indirect calorimetry. Clinical relevance: assessing nutritional needs, managing obesity, and diagnosing hyper- or hypometabolic states.

\medterm{Bile Salt Emulsification}-- Bile salts (derived from cholesterol in the liver) are amphipathic molecules with hydrophobic and hydrophilic regions that act as biological detergents, emulsifying dietary lipids into smaller droplets (micelles). This increases the surface area of lipid substrate available for pancreatic lipase action. Bile salts form mixed micelles with phospholipids, cholesterol, and fat-soluble vitamins (A, D, E, K), facilitating their transport across the unstirred water layer to the enterocyte brush border for absorption. Bile salts are reabsorbed in the terminal ileum via the apical sodium-dependent bile acid transporter (ASBT) and returned to the liver through the enterohepatic circulation, recycling 6-10 times daily.

\medterm{Blood Pressure}-- The force exerted by blood against arterial walls, measured in millimeters of mercury (mmHg). Systolic pressure (peak during ventricular ejection, normally approximately 120 mmHg) and diastolic pressure (minimum during ventricular relaxation, normally approximately 80 mmHg). Mean arterial pressure (MAP) is approximately diastolic + one-third pulse pressure (approximately 93 mmHg). MAP is the average pressure driving blood through the systemic circulation. Blood pressure depends on cardiac output and systemic vascular resistance (MAP = CO x SVR). It is regulated by short-term mechanisms (baroreceptor reflex, chemoreceptor reflex, and hormonal responses: ADH, adrenaline) and long-term mechanisms (renin-angiotensin-aldosterone system and renal body fluid regulation by the kidneys).

\medterm{Bohr Effect}-- The decrease in hemoglobin's oxygen affinity (right shift of the dissociation curve) caused by decreased pH (increased H+) and increased PCO2. In metabolically active tissues, CO2 production and lactic acid accumulation lower the local pH, causing hemoglobin to release more oxygen precisely where it is needed most. The mechanism involves protonation of hemoglobin histidine residues, which stabilize the T (tense, low-affinity) state. The Bohr effect is quantified as approximately 0.48 pH units per log10 PCO2 change and is clinically significant in conditions that alter acid-base status, enhancing oxygen delivery in exercising muscle and hypoxic tissues.

\medterm{Calciferol}-- See Vitamin D: Calcitriol.

\medterm{Calorie}-- A unit of energy measurement used in nutrition to quantify the energy content of foods and energy expenditure of the body. One kilocalorie (kcal, commonly called a Calorie with capital C) equals the energy required to raise 1 kilogram of water by 1 degree Celsius. The body derives energy from macronutrients: carbohydrates and proteins provide 4 kcal/g, fats provide 9 kcal/g, and alcohol provides 7 kcal/g. Total daily energy expenditure comprises basal metabolic rate (60--75%), thermic effect of food (10%), and physical activity (15--30%). Caloric balance determines body weight: excess calories are stored as fat; caloric deficit leads to weight loss.

\medterm{Capillary}-- The smallest blood vessel in the body, forming a microscopic network (capillary bed) that connects arterioles to venules. Capillary walls consist of a single layer of endothelial cells (frequently fenestrated) and a basement membrane, allowing exchange of oxygen, carbon dioxide, nutrients, wastes, and other substances between blood and interstitial fluid. Types: continuous (most common, in muscle, skin, lung, CNS--blood-brain barrier), fenestrated (kidney, intestine, endocrine glands, with pores for rapid filtration), and sinusoidal\/discontinuous (liver, bone marrow, spleen with large gaps allowing passage of cells and large molecules). Starling forces--hydrostatic and oncotic pressures--govern fluid movement across capillary walls.

\medterm{Cardiac Cycle}-- The sequence of mechanical and electrical events occurring during one heartbeat, lasting approximately 0.8 seconds at 75 bpm. The cycle has two main phases: systole (ventricular contraction and ejection) and diastole (ventricular relaxation and filling). Isovolumetric contraction: all valves closed, ventricles contract, pressure rises rapidly. Ventricular ejection: aortic/pulmonary valves open, blood is ejected (stroke volume approximately 70 mL). Isovolumetric relaxation: all valves closed, ventricles relax, pressure falls rapidly below atrial pressure. Ventricular filling (diastole): AV valves open, passive filling occurs (80 percent), followed by atrial contraction (20 percent, the "atrial kick"). Heart sounds: S1 (AV valve closure), S2 (semilunar valve closure), S3 (rapid ventricular filling, normal in children, abnormal in adults indicating heart failure), S4 (atrial contraction against a stiff ventricle).

\medterm{Cardiac Output}-- The volume of blood pumped by each ventricle per minute, calculated as heart rate multiplied by stroke volume (CO = HR x SV). Normal resting CO is approximately 5 liters per minute (70 mL stroke volume x 70 bpm). Stroke volume depends on preload (Frank-Starling), afterload (resistance to ejection), and contractility (inotropic state). Heart rate is modulated by the autonomic nervous system: sympathetic stimulation increases HR (positive chronotropy) via norepinephrine acting on beta-1 receptors; parasympathetic (vagal) stimulation decreases HR (negative chronotropy) via acetylcholine acting on M2 receptors. Cardiac output increases during exercise up to 4-5 times resting levels through combined increases in HR and stroke volume, with trained athletes achieving higher maximal CO.

\medterm{Cholecystokinin}-- CCK is a peptide hormone produced by I-cells in the duodenal and jejinal mucosa in response to fatty acids (long-chain fatty acids are the most potent stimuli) and amino acids in the chyme. CCK's primary actions: stimulates gallbladder contraction (promoting bile release for fat emulsification and absorption), stimulates pancreatic enzyme secretion (acinar cells), relaxes the sphincter of Oddi, slows gastric emptying (reducing the rate of acid delivery to the duodenum), and promotes satiety (acting on the hypothalamus via CCK-B receptors). CCK is also found in the brain as a neurotransmitter. Its secretion is inhibited by trypsin and chymotrypsin (negative feedback via trypsin-sensitive CCK-releasing factor). CCK analogues are used as diagnostic agents in gallbladder imaging.

\medterm{Cholesterol Synthesis}-- Occurs in all nucleated cells (primarily liver and intestine) through a complex 30+ step pathway beginning with acetyl-CoA. Key steps: three acetyl-CoA molecules condense to form HMG-CoA, which is reduced to mevalonate by HMG-CoA reductase (the rate-limiting enzyme, located in the endoplasmic reticulum). Mevalonate is converted through multiple steps to squalene, then lanosterol, then cholesterol. HMG-CoA reductase is the target of statins (competitive inhibitors), the most widely prescribed cholesterol-lowering drugs. Cholesterol is the precursor for bile salts, steroid hormones (cortisol, aldosterone, sex hormones), and vitamin D, and is a key component of cell membranes where it modulates fluidity.

\medterm{Competitive Inhibition}-- A type of enzyme inhibition where the inhibitor (I) competes with the substrate (S) for binding to the enzyme's active site. The inhibitor structurally resembles the substrate and binds reversibly to the free enzyme (E), forming an EI complex that cannot proceed to product. The apparent Km increases (lower apparent affinity for substrate) while Vmax remains unchanged (since the inhibition can be overcome by sufficiently high substrate concentrations). On a Lineweaver-Burk plot: lines intersect on the y-axis (same 1/Vmax) but have different slopes and x-intercepts (increasing apparent Km). Competitive inhibition is reversible and clinically exploited by enzyme inhibitors such as statins (HMG-CoA reductase), ACE inhibitors, and many chemotherapeutic agents that resemble natural substrates.

\medterm{Dead Space}-- The volume of ventilated air that does not participate in gas exchange. Anatomic dead space: the conducting airways (nose, pharynx, trachea, bronchi, terminal bronchioles) that warm, humidify, and filter inspired air but contain no alveoli for gas exchange. Approximately 150 mL in a 70 kg adult. Physiologic dead space includes anatomic dead space plus alveolar dead space (ventilated alveoli that are not perfused, as in pulmonary embolism). Normally, physiologic dead space equals anatomic dead space (approximately 150 mL), but in pulmonary embolism, hypotension, or lung disease, alveolar dead space increases, reducing effective alveolar ventilation. The dead space fraction (VD/VT) is measured by the Bohr equation and is clinically important in assessing ventilatory efficiency.

\medterm{Diffusion}-- The net movement of molecules from an area of higher concentration to lower concentration due to random thermal motion (Brownian motion). Fick's first law states that the rate of diffusion is proportional to the concentration gradient, surface area, and diffusion coefficient, and inversely proportional to membrane thickness: J = -P x A x (C1-C2)/d, where J is flux, P is the permeability coefficient, A is surface area, and d is thickness. Simple diffusion occurs for small, nonpolar molecules (O2, CO2, N2, steroid hormones) that can dissolve in and pass through the lipid bilayer without a transporter. Factors affecting diffusion rate include: concentration gradient, temperature, molecular size, lipid solubility, membrane surface area, and membrane thickness. In biological systems, diffusion is effective over short distances (less than 100 micrometres), which is why capillaries are narrow and alveoli are thin-walled.

\medterm{DNA Double Helix}-- The molecular structure of DNA described by Watson and Crick in 1953, consisting of two antiparallel polynucleotide chains wound around each other in a right-handed helix. The sugar-phosphate backbone is on the exterior; nitrogenous bases project inward. Complementary base pairing: adenine (A) pairs with thymine (T) via two hydrogen bonds; guanine (G) pairs with cytosine (C) via three hydrogen bonds. The helix has a major groove (22 Angstrom) and minor groove (12 Angstrom) where proteins interact with specific base sequences without unwinding the DNA. This structure enables information storage, replication, transcription, and repair, and is the foundation of molecular genetics.

\medterm{DNA Repair Mechanisms}-- Cellular processes that correct DNA damage caused by replication errors, chemicals, radiation, and reactive oxygen species. Mismatch repair (MMR): corrects base-base mismatches and insertion-deletion loops that escape proofreading during replication. MutS (MSH2/MSH6 in eukaryotes) recognizes the mismatch; MutL (MLH1/PMS2) nicks the newly synthesized strand; exonuclease removes the error; DNA polymerase and ligase fill the gap. Deficiency causes Lynch syndrome (hereditary nonpolyposis colorectal cancer, HNPCC), the most common cause of hereditary colorectal cancer. Base excision repair (BER): repairs small, non-helix-distorting lesions such as oxidised or alkylated bases. Nucleotide excision repair (NER): repairs bulky, helix-distorting lesions such as UV-induced pyrimidine dimers. Double-strand break repair: homologous recombination (HR, error-free, uses sister chromatid as template) and non-homologous end joining (NHEJ, error-prone, active in G1 phase).

\medterm{DNA Replication}-- The semi-conservative process by which DNA duplicates before cell division, producing two identical daughter molecules each containing one original and one new strand. Key steps: (1) initiation at origins of replication (oriC in bacteria, multiple origins in eukaryotes), where helicase (DnaB in bacteria, MCM complex in eukaryotes) unwinds the double helix; (2) single-strand binding proteins (SSB) stabilize unwound DNA; (3) primase (DnaG) synthesizes short RNA primers; (4) DNA polymerase III (prokaryotes) or DNA polymerase delta/epsilon (eukaryotes) extends the primer, adding nucleotides to the 3' end; (5) on the leading strand, synthesis is continuous; on the lagging strand, synthesis is discontinuous, producing Okazaki fragments; (6) RNA primers are removed by DNA polymerase I (prokaryotes) or RNase H and FEN1 (eukaryotes); (7) DNA ligase seals the nicks between fragments. Telomerase extends telomeres to prevent chromosome shortening in germ cells and stem cells.

\medterm{Electrolytes}-- Ions in body fluids that carry an electrical charge and are essential for normal physiologic function. Major electrolytes: sodium (Na+, 135--145 mEq/L, primary extracellular cation, regulates fluid balance and membrane potential), potassium (K+, 3.5--5.0 mEq/L, primary intracellular cation, critical for cardiac and neuromuscular function), chloride (Cl-, 98--106 mEq/L, major extracellular anion, follows sodium), bicarbonate (HCO3-, 22--28 mEq/L, buffer in acid-base balance), calcium (Ca2+, 8.5--10.5 mg/dL, bone structure, muscle contraction, coagulation), magnesium (Mg2+, 1.7--2.2 mg/dL, enzyme cofactor, neuromuscular stability), and phosphate (PO4 3-, 2.5--4.5 mg/dL, energy metabolism, bone structure). Electrolyte imbalances cause a wide range of symptoms and are routinely measured by basic metabolic panel or comprehensive metabolic panel.

\medterm{Electron Transport Chain}-- A series of protein complexes (I-IV) and mobile carriers embedded in the inner mitochondrial membrane that transfer electrons from NADH and FADH2 to molecular oxygen (the final electron acceptor), generating a proton gradient used for ATP synthesis. Complex I (NADH dehydrogenase): accepts electrons from NADH, pumps 4 H+ across the membrane. Complex II (succinate dehydrogenase): accepts electrons from FADH2, does not pump protons. Coenzyme Q (ubiquinone) transfers electrons from Complexes I and II to Complex III; Complex III (cytochrome bc1) transfers electrons to cytochrome c (pumps 4 H+); Complex IV (cytochrome c oxidase) transfers electrons to O2, forming H2O (pumps 2 H+). The resulting proton gradient powers ATP synthase (Complex V), which produces approximately 2.5 ATP per NADH and 1.5 ATP per FADH2. Rotenone (Complex I inhibitor), antimycin A (Complex III), and cyanide/carbon monoxide (Complex IV) are potent toxins.

\medterm{Endocytosis}-- Process by which cells internalize extracellular material by engulfing it with the cell membrane, forming intracellular vesicles. Phagocytosis ("cell eating") involves engulfing large particles (>0.5 micrometers) such as bacteria or debris, forming a phagosome that fuses with lysosomes for digestion. This is performed by macrophages, neutrophils, and dendritic cells as part of innate immunity. Pinocytosis ("cell drinking") non-specifically engulfs extracellular fluid and dissolved solutes in small vesicles (approximately 0.1 micrometres). Receptor-mediated endocytosis (clathrin-dependent) is highly specific: ligands bind to surface receptors, which cluster in clathrin-coated pits, forming coated vesicles that are uncoated and delivered to early endosomes. Examples include LDL uptake (deficiency causes familial hypercholesterolaemia), transferrin-bound iron uptake, and growth factor internalisation.

\medterm{Exocytosis}-- Process by which intracellular vesicles fuse with the plasma membrane, releasing their contents into the extracellular space and incorporating membrane proteins/lipids into the cell membrane. The SNARE complex (v-SNARE on vesicles and t-SNARE on target membrane) mediates vesicle docking and fusion. Calcium-triggered exocytosis occurs in neurons (neurotransmitter release at synapses), endocrine cells (hormone secretion), and exocrine cells (enzyme secretion). Constitutive exocytosis continuously delivers membrane components and secreted proteins (collagen, matrix proteins) in all cells. Regulated exocytosis stores secretory vesicles until a signal (typically increased cytosolic Ca2+) triggers fusion. Defects in exocytosis underlie diseases such as botulism (toxin cleaves SNARE proteins) and familial haemophagocytic lymphohistiocytosis.

\medterm{Facilitated Diffusion}-- Passive transport of molecules across cell membranes via specific transmembrane protein carriers or channels, moving substances down their concentration gradient without energy expenditure. Channel-mediated facilitated diffusion uses ion channels (voltage-gated, ligand-gated, or mechanically-gated) for rapid ion movement. Carrier-mediated facilitated diffusion uses specific carrier proteins that undergo conformational changes to transport molecules like glucose (via GLUT transporters) and amino acids. GLUT1 (erythrocytes, brain), GLUT2 (liver, pancreas), GLUT3 (neurons), GLUT4 (muscle, adipose, insulin-responsive). The rate of facilitated diffusion is limited by transporter number (saturable, following Michaelis-Menten kinetics) and is specific, saturable, and inhibitable.

\medterm{Filtration Fraction}-- The fraction of renal plasma flow that is filtered at the glomerulus, calculated as GFR divided by renal plasma flow (RPF). Normal GFR is approximately 125 mL/min and RPF is approximately 625 mL/min, yielding a filtration fraction of approximately 0.20 (20 percent). This means that approximately 20 percent of the plasma entering the glomerulus is filtered; the remaining 80 percent continues through the efferent arteriole to the peritubular capillaries. Changes in filtration fraction reflect changes in glomerular capillary hydrostatic pressure, plasma oncotic pressure, or renal blood flow. Increased filtration fraction (as in diabetic nephropathy) raises peritubular capillary oncotic pressure, enhancing reabsorption. Decreased filtration fraction suggests reduced GFR relative to RPF, as seen in prerenal azotaemia.

\medterm{Frank-Starling Mechanism}-- The relationship between ventricular end-diastolic volume (preload) and stroke volume: within physiological limits, increased venous return stretches the ventricular wall, increasing the force of contraction and stroke volume. This occurs because stretching the myocardium optimizes the overlap between actin and myosin filaments, increasing the number of active cross-bridges and enhancing calcium sensitivity. The mechanism allows the heart to automatically match its output to the volume of blood returning to it, balancing left and right ventricular outputs over successive beats. It is the basis for the clinical use of fluid resuscitation in heart failure and the relationship between central venous pressure and cardiac output on the Starling curve.

\medterm{G-Protein Coupled Receptors}-- The largest family of cell surface receptors with seven transmembrane alpha-helical domains. Upon ligand binding, the receptor undergoes a conformational change that activates an associated heterotrimeric G-protein (Galpha, Gbeta, Ggamma). The Galpha subunit exchanges GDP for GTP and dissociates from the Gbeta-gamma complex. Gs-proteins activate adenylyl cyclase (increasing cAMP), Gi-proteins inhibit adenylyl cyclase (decreasing cAMP), and Gq-proteins activate phospholipase C (producing IP3 and DAG, increasing intracellular Ca2+). The system provides signal amplification and is the target of many drugs (beta-blockers, antihistamines, opioids) and bacterial toxins (cholera toxin permanently activates Gs, pertussis toxin inactivates Gi).

\medterm{Galactosemia}-- An inborn error of galactose metabolism, most commonly caused by deficiency of galactose-1-phosphate uridylyltransferase (GALT, autosomal recessive), which converts galactose-1-phosphate to glucose-1-phosphate. Deficiency causes accumulation of galactose-1-phosphate and galactitol (from galactose reduction by aldose reductase) in tissues. Newborn screening detects elevated galactose. Symptoms appear after milk feeding (containing lactose, which is hydrolysed to glucose and galactose): vomiting, diarrhoea, failure to thrive, hypoglycaemia, hepatomegaly, jaundice, and cataracts. Treatment is a galactose-restricted diet. Long-term complications despite dietary management include intellectual disability, speech delay, and premature ovarian failure in females.

\medterm{Gastrin Secretion}-- Gastrin is a peptide hormone produced by G-cells in the gastric antrum and duodenum, stimulating hydrochloric acid secretion by parietal cells, pepsinogen secretion by chief cells, and growth of the gastric mucosa. Stimuli for release include: gastric distension, amino acids (especially phenylalanine and tryptophan) and peptides in the stomach, and vagal stimulation (via GRP, gastrin-releasing peptide). Gastrin acts on cholecystokinin-B (CCK-B) receptors on parietal cells and histamine-releasing enterochromaffin-like (ECL) cells (which in turn stimulate parietal cells via histamine H2 receptors). Gastrin release is inhibited by low antral pH (negative feedback via somatostatin from D-cells). Zollinger-Ellison syndrome (gastrinoma) causes severe peptic ulcer disease from unregulated gastrin secretion.

\medterm{Gaucher Disease}-- The most common lysosomal storage disease, caused by deficiency of glucocerebrosidase (glucosylceramidase, beta-glucocerebrosidase, GBA), an autosomal recessive disorder leading to accumulation of glucocerebroside (glucosylceramide) in macrophages of the reticuloendothelial system. These lipid-laden macrophages are called "Gaucher cells" with a characteristic "crinkled tissue paper" cytoplasmic appearance. Type 1 (non-neuronopathic, most common, approximately 90 percent of cases): hepatosplenomegaly, anaemia, thrombocytopenia, bone pain crises, and Erlenmeyer flask deformity of the distal femurs. Type 2 (acute neuronopathic, infantile): severe neurodegeneration with brainstem involvement and death by age 2. Type 3 (chronic neuronopathic): later onset with myoclonus, seizures, and progressive neurological deterioration.

\medterm{Gene Therapy}-- The therapeutic delivery of nucleic acids into cells to treat or prevent disease by altering gene expression. Two main approaches: (1) gene replacement/supplementation (adding a functional copy of a defective gene), and (2) gene editing (correcting the defective gene in situ). Viral vectors: adeno-associated virus (AAV, non-integrating, low immunogenicity, limited packaging capacity), lentivirus (integrates into the host genome, effective for dividing and non-dividing cells), and adenovirus (episomal, highly immunogenic, large packaging capacity). Non-viral methods include lipid nanoparticles (mRNA vaccines, siRNA delivery) and electroporation. Major challenges include immunogenicity, off-target effects, delivery efficiency, and long-term expression durability.

\medterm{Glomerular Filtration Rate}-- The volume of plasma filtered through the glomerular capillaries per unit time, normally approximately 125 mL/min or approximately 180 L/day. GFR depends on Starling forces across the glomerular capillary wall: glomerular capillary hydrostatic pressure (approximately 55 mmHg, favoring filtration) minus Bowman's capsule hydrostatic pressure (approximately 15 mmHg) minus glomerular capillary oncotic pressure (approximately 30 mmHg, opposing filtration), yielding a net filtration pressure of approximately 10 mmHg. GFR is autoregulated between MAP of approximately 80-180 mmHg by the myogenic response (afferent arteriolar constriction in response to increased pressure) and tubuloglomerular feedback (macula densa sensing NaCl delivery). Clinical GFR measurement uses creatinine clearance or cystatin C. Reduced GFR defines acute kidney injury and chronic kidney disease staging.

\medterm{Gluconeogenesis}-- The synthesis of glucose from non-carbohydrate precursors (lactate, glycerol, glucogenic amino acids, and propionate) primarily in the liver (90 percent) and kidneys (10 percent). It uses most glycolytic enzymes in reverse but must bypass three irreversible glycolytic reactions using different enzymes: (1) pyruvate carboxylase converts pyruvate to oxaloacetate, then PEPCK (phosphoenolpyruvate carboxykinase) converts oxaloacetate to phosphoenolpyruvate; (2) fructose-1,6-bisphosphatase converts fructose-1,6-bisphosphate to fructose-6-phosphate; (3) glucose-6-phosphatase (present in liver and kidney but not muscle) converts glucose-6-phosphate to free glucose for release into the blood. Gluconeogenesis is energy-intensive (6 ATP equivalents per glucose) and is stimulated by glucagon, cortisol, and catecholamines, and inhibited by insulin and ethanol.

\medterm{Glycogen Storage Diseases}-- A group of inherited disorders caused by deficiencies in enzymes involved in glycogen metabolism, affecting the liver, muscle, or both. Type I (von Gierke disease): glucose-6-phosphatase deficiency, causing severe fasting hypoglycemia, hepatomegaly (massive glycogen and fat accumulation), lactic acidosis, hyperuricemia, and hyperlipidemia. Type II (Pompe disease): acid alpha-glucosidase (lysosomal acid maltase) deficiency, causing accumulation of glycogen in lysosomes, leading to cardiomegaly, hypertrophic cardiomyopathy, and muscle weakness; infantile form is fatal without enzyme replacement therapy. Type III (Cori disease): debranching enzyme deficiency, milder hypoglycaemia and hepatomegaly. Type V (McArdle disease): muscle glycogen phosphorylase deficiency, causing exercise intolerance, myoglobinuria, and second wind phenomenon.

\medterm{Glycogenesis}-- The synthesis of glycogen from glucose, primarily in liver (for blood glucose maintenance) and skeletal muscle (for local energy use). The process involves two key enzymes: glycogen synthase (adding glucose residues from UDP-glucose to the non-reducing ends of glycogen via alpha-1,4-glycosidic bonds) and the branching enzyme (creating alpha-1,6-glycosidic bonds every 8-12 residues, creating branch points). Glycogen synthase is the rate-limiting enzyme, regulated by phosphorylation (inactivated by PKA, activated by dephosphorylation via protein phosphatase 1) and allosteric activation by glucose-6-phosphate. Insulin promotes glycogenesis by activating protein phosphatase 1 and inactivating glycogen synthase kinase 3 (GSK-3). Glucagon and adrenaline inhibit glycogenesis via cAMP-mediated phosphorylation.

\medterm{Glycogenolysis}-- The breakdown of glycogen to glucose, primarily in the liver (to maintain blood glucose) and muscle (for local energy use). Key enzyme: glycogen phosphorylase, which cleaves alpha-1,4-glycosidic bonds by phosphorylating glucose residues to produce glucose-1-phosphate. The debranching enzyme handles branch points (alpha-1,6 bonds). In the liver, glucose-1-phosphate is converted to glucose-6-phosphate by phosphoglucomutase, then to free glucose by glucose-6-phosphatase (released into blood). In muscle, glucose-6-phosphate enters glycolysis directly (cannot be released as free glucose). Glycogenolysis is activated by glucagon (liver), adrenaline (liver and muscle), and AMP (muscle, during exercise), and inhibited by insulin and glucose.

\medterm{Glycolysis}-- The metabolic pathway converting glucose to pyruvate, occurring in the cytoplasm of all cells. Ten enzymatic steps divided into two phases: the energy investment phase (steps 1-5, consuming 2 ATP) and the energy payoff phase (steps 6-10, producing 4 ATP and 2 NADH). Key regulatory enzymes: hexokinase/glucokinase (step 1, phosphorylating glucose), phosphofructokinase-1 (PFK-1, step 3, the rate-limiting step, allosterically activated by AMP, fructose-2,6-bisphosphate, and inhibited by ATP and citrate). The net yield per glucose is 2 ATP and 2 NADH (under aerobic conditions); under anaerobic conditions, NAD+ is regenerated by lactate dehydrogenase reducing pyruvate to lactate, allowing continued glycolysis but with net 2 ATP per glucose.

\medterm{Graded Potentials}-- Local changes in membrane potential that vary in magnitude proportional to stimulus strength, unlike all-or-none action potentials. They include receptor potentials (in sensory receptors), postsynaptic potentials (EPSPs and IPSPs at synapses), and pacemaker potentials (in cardiac/specialized cells). Graded potentials decay with distance from their source (decremental conduction) because they spread passively through the cytoplasm. They can summate temporally (rapid successive stimuli) and spatially (simultaneous stimuli from multiple sources), integrating signals at the axon hillock. If the summed depolarisation reaches threshold (approximately -55 mV), an action potential is triggered. Graded potentials are essential for sensory transduction and synaptic communication and are the basis for the all-or-none law of neural signalling.

\medterm{Haldane Effect}-- The decrease in hemoglobin's CO2 carrying capacity when oxygen binds to it, facilitating CO2 release at the lungs. Deoxygenated hemoglobin (deoxyhemoglobin) is a better buffer for H+ ions and forms more carbaminohemoglobin than oxyhemoglobin. In the tissues: oxygen unloading increases hemoglobin's ability to bind CO2 (as carbaminohemoglobin) and buffer H+ (from CO2 hydration to carbonic acid), enhancing CO2 loading. In the lungs: oxygen binding displaces CO2 from hemoglobin and reduces H+ buffer capacity, promoting CO2 release into the alveoli. The Haldane effect explains that approximately 20 percent of CO2 transport is via carbaminohemoglobin. Together with the Bohr effect, it optimizes O2 delivery and CO2 removal, a concept known as the "Haldane-Bohr coupling" that is essential for efficient gas exchange.

\medterm{Homeostasis}-- The maintenance of a stable internal environment despite external changes, achieved through physiological regulation and feedback mechanisms. The concept was introduced by Walter Cannon (1929) building on Claude Bernard's milieu intérieur. Homeostasis involves: (1) Set point: the ideal or normal value for a physiological variable (e.g., body temperature 37 C, blood pH 7.4, blood glucose 4-6 mmol/L). (2) Sensors (receptors): detect deviations from the set point (e.g., thermoreceptors in the hypothalamus for temperature, chemoreceptors for pH and O2/CO2, baroreceptors for blood pressure). (3) Integrating centre: compares input to the set point and coordinates a response (e.g., hypothalamus, brainstem, endocrine glands). (4) Effectors: muscles, glands, or organs that produce the corrective response (e.g., sweat glands and shivering for temperature; insulin and glucagon for glucose; heart, vessels, and kidneys for BP). Feedback mechanisms: (1) Negative feedback (most common): reverses the direction of change (e.g., rising blood glucose stimulates insulin release, which lowers glucose; blood pressure increase triggers baroreceptor-mediated vasodilation and bradycardia). (2) Positive feedback: amplifies the initial change, driving a process to completion (e.g., oxytocin release during childbirth intensifies uterine contractions; sodium channel opening during action potential depolarisation; blood clotting cascade). Disorders of homeostasis: failure of temperature regulation (hyperthermia, hypothermia), pH (acidosis, alkalosis), glucose (diabetes, hypoglycaemia), blood pressure (hypertension, hypotension), and water/electrolyte balance (dehydration, hyponatraemia, hyperkalaemia).

\medterm{Hypoxic Pulmonary Vasoconstriction}-- The unique response of pulmonary arterioles to low alveolar oxygen tension (PO2 less than approximately 60 mmHg), causing vasoconstriction. Unlike systemic vessels that dilate in response to hypoxia, pulmonary vessels constrict, redirecting blood flow away from poorly ventilated alveoli toward better-ventilated regions, optimizing ventilation-perfusion (V/Q) matching. The mechanism involves inhibition of voltage-gated K+ channels in pulmonary artery smooth muscle cells, causing membrane depolarization, Ca2+ influx via L-type Ca2+ channels, and vasoconstriction. Chronic global hypoxia (as in COPD, high altitude) causes sustained vasoconstriction leading to pulmonary hypertension, right ventricular hypertrophy (cor pulmonale), and vascular remodelling. This response is unique to the pulmonary circulation; systemic vessels dilate in hypoxia.

\medterm{JAK-STAT Pathway}-- A rapid signaling pathway for cytokines and growth factors (interferons, interleukins, erythropoietin, growth hormone) that lack intrinsic enzymatic activity. Ligand binding to the cytokine receptor causes receptor dimerization and activation of associated Janus kinases (JAK1, JAK2, JAK3, TYK2) by trans-phosphorylation. Activated JAKs phosphorylate tyrosine residues on the receptor cytoplasmic domain, creating docking sites for STAT (signal transducers and activators of transcription) proteins, which are phosphorylated by JAKs, form homo- or heterodimers, and translocate to the nucleus where they bind gamma-interferon activation site (GAS) or interferon-stimulated response element (ISRE) sequences to regulate gene transcription. The pathway is negatively regulated by SOCS (suppressor of cytokine signalling) proteins, PTPs (protein tyrosine phosphatases), and PIAS (protein inhibitor of activated STAT).

\medterm{Ketogenesis}-- The hepatic synthesis of ketone bodies (acetoacetate, beta-hydroxybutyrate, and acetone) from acetyl-CoA during periods of prolonged fasting, starvation, uncontrolled diabetes, or very low carbohydrate diets. When oxaloacetate is depleted for gluconeogenesis, accumulated acetyl-CoA cannot enter the TCA cycle and is diverted to ketogenesis. The process occurs in the mitochondrial matrix: two acetyl-CoA molecules condense to form acetoacetyl-CoA (by thiolase), which combines with another acetyl-CoA to form HMG-CoA by HMG-CoA synthase. HMG-CoA lyase cleaves HMG-CoA to acetoacetate and acetyl-CoA. Acetoacetate is reduced to beta-hydroxybutyrate by beta-hydroxybutyrate dehydrogenase (NaDH-dependent) or spontaneously decarboxylated to acetone. Ketone bodies are important alternative fuels for the brain during fasting and starvation, providing up to 70 percent of cerebral energy needs.

\medterm{Lipoprotein Metabolism}-- Lipoproteins transport hydrophobic lipids (triglycerides, cholesterol, cholesterol esters, fat-soluble vitamins) in the aqueous blood. Four main classes: chylomicrons (from intestine, largest, transport dietary triglycerides; ApoB-48; acted on by lipoprotein lipase), VLDL (from liver, transport endogenous triglycerides; ApoB-100; converted to IDL), LDL (from VLDL metabolism, carry cholesterol to tissues; ApoB-100; the "bad" cholesterol; taken up by LDL receptors), and HDL (from liver/intestine, reverse cholesterol transport, ApoA-I, the "good" cholesterol, protective against atherosclerosis). Lipoprotein lipase (LPL) hydrolyses triglycerides in chylomicrons and VLDL, releasing free fatty acids for tissue uptake. HDL mediates reverse cholesterol transport, delivering excess cholesterol from peripheral cells to the liver for excretion or recycling.

\medterm{Membrane Potential}-- The electrical potential difference across a cell membrane, approximately -70 mV in most neurons (resting membrane potential). Established by the unequal distribution of ions (primarily K+, Na+, Cl-) across the membrane and selective permeability. The Nernst equation calculates the equilibrium potential for a single ion: E = (RT/zF) ln([ion outside]/[ion inside]). At body temperature for a monovalent cation, this simplifies to E = 61.5 log([out]/[in]) mV. The Goldman equation accounts for multiple ion permeabilities to calculate the resting membrane potential: Vm = 61.5 log((PK[K+]o + PNa[Na+]o + PCl[Cl-]i)/(PK[K+]i + PNa[Na+]i + PCl[Cl-]o)). The resting membrane is most permeable to K+ (via leak channels), so Vm is close to EK (approximately -94 mV for K+). Changes in ion permeability (e.g., Na+ channels opening during depolarisation) drive electrical signalling.

\medterm{Michaelis-Menten Kinetics}-- The mathematical model describing the rate of enzyme-catalyzed reactions as a function of substrate concentration. The Michaelis-Menten equation: v = (Vmax x [S]) / (Km + [S]), where v is the reaction velocity, Vmax is the maximum velocity (when all enzyme is saturated with substrate), [S] is the substrate concentration, and Km is the Michaelis constant (the substrate concentration at which v = Vmax/2). Km reflects the enzyme's affinity for its substrate: a low Km indicates high affinity (the enzyme reaches half-maximal velocity at a low substrate concentration). The turnover number (kcat = Vmax/[Et]) describes the catalytic efficiency; the ratio kcat/Km is the best measure of enzyme specificity, approaching the diffusion limit for perfect enzymes. Lineweaver-Burk (double reciprocal) plots linearise the equation: 1/v = Km/Vmax x 1/[S] + 1/Vmax.

\medterm{Migrating Motor Complex}-- Cyclical patterns of intense gastrointestinal motility occurring during fasting (interdigestive period), sweeping residual food, bacteria, and debris from the stomach through the small intestine to the colon. Each cycle lasts approximately 90-120 minutes and has four phases: Phase I (quiescence, approximately 45-60 minutes), Phase II (irregular contractions, approximately 30 minutes), Phase III (regular high-amplitude contractions at maximum frequency, approximately 5-15 minutes, the "housekeeper" phase that propels luminal contents distally), and Phase IV (declining activity returning to Phase I). The MMC is regulated by the migrating motor complex, motilin (released during fasting), and the enteric nervous system. Disruption of the MMC is associated with small intestinal bacterial overgrowth (SIBO).

\medterm{Na+/K+ ATPase}-- The sodium-potassium pump is a P-type ATPase found in the plasma membrane of all animal cells, directly consuming approximately 20-30 percent of cellular ATP. It pumps 3 Na+ ions out of the cell and 2 K+ ions in per ATP hydrolyzed, establishing steep electrochemical gradients: high Na+ outside, high K+ inside. This gradient is essential for maintaining the resting membrane potential (approximately -70 mV), driving secondary active transport (e.g., glucose and amino acid absorption in the intestine and kidney, Na+/Ca2+ exchange), and regulating cell volume. The pump is inhibited by cardiac glycosides such as ouabain and digoxin (used in heart failure to increase intracellular Ca2+ and contractility).

\medterm{Nephron Function}-- The nephron is the functional unit of the kidney, with approximately 1 million per kidney. Each nephron consists of a renal corpuscle (glomerulus and Bowman's capsule) and a renal tubule (proximal convoluted tubule, loop of Henle, distal convoluted tubule, and collecting duct). The nephron performs three primary functions: filtration (at the glomerulus, producing approximately 180 L/day of ultrafiltrate), reabsorption ( reclaiming approximately 99 percent of filtered water and solutes), and secretion (transporting substances from the blood into the tubular lumen for elimination). The loop of Henle establishes the medullary osmotic gradient for urine concentration; the collecting duct performs final regulation of water and electrolyte balance under the control of ADH and aldosterone.

\medterm{Noncompetitive Inhibition}-- A type of enzyme inhibition where the inhibitor binds at a site distinct from the active site (allosteric site), and can bind to both the free enzyme (E) and the enzyme-substrate complex (ES) with equal affinity. The inhibitor does not compete with the substrate and cannot be overcome by increasing substrate concentration. The apparent Vmax decreases (fewer functional enzyme molecules) while Km remains unchanged (substrate affinity is unaffected since the inhibitor doesn't compete for the active site). On a Lineweaver-Burk plot: lines intersect on the x-axis (same -1/Km) with different slopes and y-intercepts. Noncompetitive inhibition is less common than competitive inhibition in pharmacology but is important in metal ion toxicity and certain pesticide actions.

\medterm{Osmosis}-- The net movement of water molecules across a semipermeable membrane from a region of lower solute concentration (higher water concentration) to a region of higher solute concentration (lower water concentration). Driven by the osmotic pressure difference, which is proportional to the number of solute particles. Osmolarity is measured in osmoles per liter (mOsm/L); normal plasma osmolarity is approximately 285-295 mOsm/L. Tonicity describes the effect of a solution on cell volume: hypotonic solutions cause cell swelling (osmotic lysis), hypertonic solutions cause cell shrinkage (crenation). The concept is clinically critical for intravenous fluid therapy: normal saline (0.9 percent NaCl, approximately 308 mOsm/L) is isotonic and does not cause red blood cell lysis; hypotonic fluids (0.45 percent NaCl) may cause haemolysis if infused rapidly.

\medterm{Oxidative Phosphorylation}-- The process by which ATP is synthesized using the energy stored in the proton gradient (proton motive force) across the inner mitochondrial membrane, established by the electron transport chain. ATP synthase (F1F0-ATPase) is a molecular turbine with two components: F0 (membrane-spanning proton channel) and F1 (catalytic head projecting into the matrix). Protons flow through F0 down their electrochemical gradient, causing the gamma subunit to rotate relative to the alpha3beta3 hexamer of F1, driving conformational changes that synthesise ATP from ADP and Pi. The chemiosmotic theory (Peter Mitchell, 1978 Nobel Prize) explains how the proton motive force (proton gradient + membrane potential) couples electron transport to ATP synthesis. Uncouplers (e.g., DNP) dissipate the proton gradient, generating heat instead of ATP.

\medterm{Oxygen-Hemoglobin Dissociation Curve}-- The sigmoidal relationship between partial pressure of oxygen (PO2) and hemoglobin oxygen saturation (SO2). The curve reflects hemoglobin's cooperative binding: oxygen binding to one subunit increases affinity for subsequent oxygen molecules. At PO2 of 100 mmHg (alveolar), hemoglobin is approximately 98 percent saturated. At PO2 of 40 mmHg (mixed venous), saturation is approximately 75 percent. The steep portion of the curve (PO2 20-60 mmHg) allows large oxygen unloading with small PO2 decreases in tissue capillaries. The curve is shifted rightward (decreased affinity) by increased temperature, decreased pH (Bohr effect), increased PCO2, and increased 2,3-BPG, enhancing oxygen unloading. Leftward shifts (increased affinity) occur with decreased temperature, increased pH, and decreased 2,3-BPG (e.g., stored blood), reducing oxygen unloading. Foetal haemoglobin (HbF) has a left-shifted curve for placental oxygen transfer.

\medterm{Pentose Phosphate Pathway}-- A cytoplasmic glucose oxidation pathway running parallel to glycolysis, primarily in the liver, adipose tissue, adrenal cortex, and red blood cells. It has two phases: the oxidative phase (producing NADPH and ribose-5-phosphate) and the non-oxidative phase (interconverting sugar phosphates). Key enzyme: glucose-6-phosphate dehydrogenase (G6PD), the rate-limiting enzyme, producing NADPH. NADPH is essential for: fatty acid and steroid biosynthesis, maintaining reduced glutathione (protecting red blood cells from oxidative damage and preventing haemolysis), and cytochrome P450 detoxification. The non-oxidative phase recycles ribose-5-phosphate to glycolytic intermediates. G6PD deficiency (X-linked, common in Mediterranean, African, and Asian populations) causes haemolytic anaemia triggered by oxidative stress from drugs (primaguine, sulphonamides) or fava beans.

\medterm{Peristalsis}-- Rhythmic sequential contractions of longitudinal and circular smooth muscle that propel gastrointestinal contents in an aboral (oral to anal) direction. The peristaltic reflex involves stretch activation of the intestinal wall: proximal to the stretch, longitudinal muscle contracts and circular muscle relaxes (ascending contraction); distally, circular muscle contracts and longitudinal muscle relaxes (descending relaxation). This is mediated by the myenteric (Auerbach's) plexus: excitatory neurotransmitters (acetylcholine, substance P) cause proximal contraction, and inhibitory neurotransmitters (VIP, nitric oxide) cause distal relaxation. Peristalsis propels chyme at a rate of approximately 1 cm/min in the small intestine, 5-10 cm/min in the oesophagus, and is slower in the colon. Disorders include achalasia (failure of LES relaxation), intestinal pseudo-obstruction, and Hirschsprung disease.

\medterm{Phenylketonuria}-- The most common inborn error of amino acid metabolism, caused by deficiency of phenylalanine hydroxylase (PAH, autosomal recessive), which converts phenylalanine to tyrosine. The deficiency leads to accumulation of phenylalanine and its transamination product phenylpyruvate in the blood and brain. Newborn screening (Guthrie test or tandem mass spectrometry) detects elevated phenylalanine levels (greater than 120 micromol/L). Untreated PKU causes severe intellectual disability, seizures, behavioural problems, eczema, microcephaly, and a musty odour. Newborn screening and early dietary restriction (low phenylalanine diet with special formula) prevents neurological damage. Maternal PKU (uncontrolled hyperphenylalaninaemia during pregnancy) causes fetal intellectual disability, microcephaly, and congenital heart defects.

\medterm{Phospholipase C Pathway}-- A GPCR signaling pathway activated primarily by Gq-protein coupled receptors (e.g., alpha-1 adrenergic, muscarinic M1/M3, vasopressin V1, angiotensin II AT1 receptors). When the Gq-alpha subunit is activated, it stimulates phospholipase C-beta (PLC-beta), which cleaves the membrane phospholipid phosphatidylinositol 4,5-bisphosphate (PIP2) into two second messengers: inositol 1,4,5-trisphosphate (IP3) and diacylglycerol (DAG). IP3 is a water-soluble molecule that diffuses through the cytoplasm and binds to IP3 receptors on the endoplasmic reticulum, triggering Ca2+ release into the cytoplasm. DAG remains in the membrane and activates protein kinase C (PKC), which phosphorylates downstream target proteins. The Ca2+ signal activates calmodulin-dependent kinases (CaMK), and the pathway mediates smooth muscle contraction, neurotransmitter release, hormone secretion, and cell proliferation.

\medterm{Phytomenadione}-- See Vitamin K: Phylloquinone.

\medterm{Protein Folding}-- The process by which newly synthesized polypeptide chains acquire their functional three-dimensional structures. The amino acid sequence (primary structure) determines the folding pathway. Key structural levels: primary (amino acid sequence), secondary (alpha-helices and beta-sheets formed by hydrogen bonding), tertiary (overall 3D fold of a single polypeptide), and quaternary (arrangement of multiple subunits). Chaperone proteins assist proper folding: Hsp70 (heat shock protein 70) binds exposed hydrophobic regions of nascent polypeptides; Hsp60 (chaperonin, GroEL/GroES in bacteria) provides an isolated compartment for folding. Misfolded proteins are targeted for degradation by the ubiquitin-proteasome system. Protein misfolding causes conformational diseases including Alzheimer's (amyloid-beta), Parkinson's (alpha-synuclein), prion diseases (Creutzfeldt-Jakob), and cystic fibrosis (CFTR misfolding).

\medterm{Proximal Tubule Reabsorption}-- The proximal convoluted tubule (PCT) reabsorbs approximately 65 percent of filtered sodium, water, glucose, amino acids, bicarbonate, and phosphate. Sodium reabsorption is the driving force for most solute and water transport: the basolateral Na+/K+ ATPase creates a low intracellular Na+ concentration, allowing Na+ entry from the tubular lumen via secondary active transporters (Na+/glucose, Na+/amino acid cotransporters on the apical membrane; Na+/H+ exchanger). Water follows sodium osmotically via aquaporins (AQP1) in the proximal tubule. Glucose and amino acids are completely reabsorbed in the early PCT via sodium-coupled transporters (SGLT2 for glucose, SGLT1 for glucose and amino acids) until their transport maximum (Tm) is exceeded. Bicarbonate is reabsorbed via H+ secretion (Na+/H+ exchanger) and carbonic anhydrase in the brush border. Phosphate reabsorption is regulated by PTH (decreases reabsorption, increases phosphate excretion).

\medterm{Pulmonary Surfactant}-- A phospholipid-protein mixture produced by type II alveolar pneumocytes that reduces surface tension at the air-liquid interface within alveoli. The primary component is dipalmitoylphosphatidylcholine (DPPC, lecithin), which forms a monolayer that reduces surface tension to near zero during expiration. Surfactant prevents alveolar collapse (atelectasis) at end-expiration by reducing the collapsing force, increases lung compliance, and stabilizes alveoli of different sizes (preventing small alveoli from collapsing into larger ones via the LaPlace relationship). Surfactant is produced from approximately week 24 of gestation, and its deficiency in premature infants causes respiratory distress syndrome (RDS) with atelectasis, hypoxaemia, and the need for surfactant replacement therapy and mechanical ventilation.

\medterm{Pyruvate}-- The three-carbon alpha-keto acid produced at the end of glycolysis from phosphoenolpyruvate by pyruvate kinase. Under aerobic conditions, pyruvate enters the mitochondrial matrix via the mitochondrial pyruvate carrier (MPC) and is converted to acetyl-CoA by the pyruvate dehydrogenase complex (PDC), an irreversible reaction that also produces CO2 and NADH. The PDC requires five cofactors derived from vitamins: thiamine pyrophosphate (B1), FAD (B2), NAD+ (B3), CoA (from pantothenic acid, B5), and lipoic acid. Under anaerobic conditions, pyruvate is reduced to lactate by lactate dehydrogenase, regenerating NAD+ for continued glycolysis. Pyruvate is the branch point between aerobic (TCA cycle) and anaerobic (lactate) metabolism, and its fate is determined by oxygen availability and tissue metabolic demands.

\medterm{Receptor Tyrosine Kinases}-- Single-pass transmembrane receptors with intrinsic tyrosine kinase activity in their intracellular domain. Ligand binding (e.g., insulin, EGF, PDGF, FGF) induces receptor dimerization and autophosphorylation of tyrosine residues. Phosphorylated tyrosines serve as docking sites for intracellular signaling proteins containing SH2 domains, activating downstream cascades: the Ras-MAP kinase pathway (cell proliferation and differentiation), the PI3-kinase/Akt pathway (cell survival and metabolism), and the PLC-gamma pathway (Ca2+ signalling). RTKs are frequently mutated or overexpressed in cancer (e.g., HER2 in breast cancer, EGFR in lung cancer), making them important therapeutic targets for tyrosine kinase inhibitors (imatinib, gefitinib, trastuzumab).

\medterm{Renal Clearance}-- The volume of plasma completely cleared of a substance by the kidneys per unit time, measured in mL/min. Clearance of any substance X is calculated as: CX = (Ux x V)/Px, where Ux is urine concentration, V is urine flow rate, and Px is plasma concentration. Inulin clearance equals GFR (approximately 125 mL/min) because inulin is freely filtered but neither reabsorbed nor secreted. Para-aminohippuric acid (PAH) clearance approximates renal plasma flow (approximately 625 mL/min) because it is both filtered and completely secreted by the proximal tubule. Creatinine clearance is the most common clinical GFR estimate; cystatin C is a more accurate alternative unaffected by muscle mass. The concept is essential for determining drug dosing in renal impairment.

\medterm{Renin-Angiotensin-Aldosterone System}-- The RAAS is a hormonal cascade regulating blood pressure, fluid volume, and electrolyte balance. Stimulus: decreased renal perfusion pressure (detected by baroreceptors in the afferent arteriole), decreased sodium delivery to the macula densa, or increased sympathetic activity. Renin (from juxtaglomerular cells) cleaves angiotensinogen (from the liver) to angiotensin I. Angiotensin-converting enzyme (ACE, primarily in pulmonary endothelium) converts angiotensin I to angiotensin II. Angiotensin II is a potent vasoconstrictor, stimulates aldosterone secretion (Na+ and water retention), increases ADH release, enhances thirst, increases sympathetic activity, and promotes cardiac and vascular remodelling. The RAAS is a target of ACE inhibitors, angiotensin receptor blockers (ARBs), direct renin inhibitors, and mineralocorticoid receptor antagonists for treating hypertension and heart failure.

\medterm{Retinol}-- See Vitamin A: Retinal.

\medterm{Salivation}-- The secretion of saliva by the salivary glands. Salivation is controlled by the autonomic nervous system: parasympathetic stimulation (CN VII—submandibular and sublingual glands; CN IX—parotid gland) produces profuse, watery saliva via acetylcholine acting on muscarinic M\(_3\) receptors; sympathetic stimulation produces scanty, viscous saliva via noradrenaline acting on \(\beta\)-adrenergic receptors. Salivation is triggered by: sight, smell, or thought of food (cephalic phase), mechanical stimulation of the mouth (food in the oral cavity), and the presence of acid or chemical irritants. The salivary centre is located in the medulla oblongata (superior and inferior salivatory nuclei). Daily volume: 1--1.5 L. Hypersalivation (sialorrhoea) can be caused by oral inflammation, teething, pregnancy, neurological disorders (Parkinson disease, motor neuron disease, cerebral palsy, rabies), gastro-oesophageal reflux, and medications (clozapine, anticholinesterases). Hyposalivation (xerostomia/dry mouth): Sjögren syndrome, radiation therapy to the head and neck, medications (anticholinergics, antidepressants, antihistamines, diuretics), diabetes, and dehydration.

\medterm{Saltatory Conduction}-- The rapid propagation of action potentials along myelinated nerve fibers, where the action potential appears to "jump" from one Node of Ranvier to the next. Myelin (produced by Schwann cells in the PNS and oligodendrocytes in the CNS) insulates the axon, preventing ion flow across the internodal membrane. Action potentials are generated only at the nodes (approximately 1 mm apart), where voltage-gated Na+ channels are concentrated. The local current from one node depolarises the next node to threshold, causing saltatory conduction at speeds up to 50 m/s in myelinated fibres (compared to approximately 0.5 m/s in unmyelinated fibres). Saltatory conduction is energy-efficient, as Na+/K+ ATPase activity is required only at the nodes. Demyelinating diseases such as multiple sclerosis (CNS) and Guillain-Barré syndrome (PNS) slow or block conduction, causing neurological deficits.

\medterm{Second Messenger Systems}-- Intracellular signaling molecules that relay extracellular signals (first messengers: hormones, neurotransmitters, growth factors) to cellular responses. Major second messengers: cAMP (produced by adenylyl cyclase from ATP, activated by Gs-GPCRs; activates PKA; degraded by phosphodiesterase; example: epinephrine binding beta-1 receptors in the heart increasing heart rate and contractility). cGMP (produced by guanylyl cyclase; activated by nitric oxide in smooth muscle causing vasodilation; also activated by natriuretic peptides ANP/BNP). IP3 (produced by phospholipase C from PIP2, releases Ca2+ from ER). DAG (produced with IP3, activates PKC). Ca2+ itself acts as a second messenger, activating calmodulin and CaMK. Second messenger systems provide signal amplification: a single activated receptor can activate many G-proteins, each activating an enzyme that produces many second messenger molecules.

\medterm{Second Messengers}-- Intracellular signaling molecules generated in response to extracellular signals (first messengers like hormones or neurotransmitters). cAMP (cyclic adenosine monophosphate): produced from ATP by adenylyl cyclase, activated by Gs-protein coupled receptors, activates protein kinase A (PKA); degraded by phosphodiesterase. IP3 (inositol 1,4,5-trisphosphate): produced from PIP2 by phospholipase C, triggers Ca2+ release from the endoplasmic reticulum. DAG (diacylglycerol): remains in the membrane along with IP3, and together they activate protein kinase C (PKC). Ca2+ is a versatile second messenger that binds calmodulin, activating CaMKII (in neurons, critical for long-term potentiation and memory formation). The second messenger concept is central to understanding hormone action, drug effects, and intracellular communication.

\medterm{Secretin}-- A peptide hormone produced by S-cells in the duodenal and jejunal mucosa in response to acidic chyme (pH less than 4.5) entering the duodenum from the stomach. Secretin's primary action is to stimulate the pancreatic ductal cells to secrete a bicarbonate-rich fluid, neutralizing duodenal acid and raising the pH to approximately 7-8, which optimizes pancreatic enzyme activity. Secretin also inhibits gastric acid secretion and gastric motility (slowing gastric emptying), stimulates bile secretion and gallbladder contraction, and enhances the effects of cholecystokinin (CCK). Secretin release is inhibited by H+ buffering in the duodenum (negative feedback). Secretin was the first hormone discovered (Bayliss and Starling, 1902), establishing the concept of hormonal regulation.

\medterm{Segmentation}-- A pattern of intestinal motility involving localized, rhythmic contractions of circular muscle that segment the intestinal contents, mixing them with digestive secretions and increasing contact with the absorptive mucosal surface. Unlike peristalsis, segmentation does not propel contents distally. It occurs predominantly in the fed state in the small intestine, at a frequency determined by the slow wave rhythm of the intestine (approximately 12 per minute in the duodenum, 8 per minute in the ileum). The contractions are mediated by the enteric nervous system and the interstitial cells of Cajal (pacemaker cells). Segmentation is essential for efficient digestion and absorption; loss of segmentation occurs in intestinal obstruction and ileus.

\medterm{Sensory Nerves}-- Nerves that transmit sensory information from peripheral receptors to the central nervous system. Sensory (afferent) neurons have their cell bodies in the dorsal root ganglia (spinal nerves) or cranial nerve ganglia (cranial nerves). Sensory receptors: mechanoreceptors (touch, pressure, vibration, stretch—Pacinian corpuscles, Meissner corpuscles, Merkel discs, Ruffini endings), thermoreceptors (cold—A\(\delta\) fibres; warm—unmyelinated C fibres), nociceptors (pain—A\(\delta\) fibres: fast, sharp, localised; C fibres: slow, dull, diffuse), chemoreceptors (taste, smell, O\(_2\), CO\(_2\), pH), and photoreceptors (rods and cones in the retina). Sensory pathways: spinothalamic tract (pain, temperature, crude touch—crosses in spinal cord), dorsal columns/medial lemniscus (fine touch, vibration, proprioception—crosses in medulla), and trigeminothalamic tract (facial sensation). Dermatomes: areas of skin supplied by a single spinal nerve (C4—shoulder, C6—thumb, C8—little finger, T4—nipple, T10—umbilicus, L4—medial foot, S1—lateral foot). Sensory deficits help localise neurological lesions at the level of the peripheral nerve, plexus, nerve root, spinal cord, brainstem, thalamus, or cortex.

\medterm{Sympathetic Nervous System}-- The division of the autonomic nervous system responsible for the 'fight or flight' response, preparing the body for intense physical activity or stress. Sympathetic preganglionic neurons originate in the intermediolateral column of the spinal cord grey matter (T1--L2), making it the thoracolumbar system. Preganglionic fibres are short and synapse in the paravertebral sympathetic chain ganglia (bilaterally, 22 ganglia) or prevertebral ganglia (coeliac, superior mesenteric, inferior mesenteric, aorticorenal). The neurotransmitter of preganglionic fibres is acetylcholine (nicotinic receptors); postganglionic fibres release noradrenaline (norepinephrine) at target organs (adrenergic \(\alpha\) and \(\beta\) receptors), except sweat glands and some blood vessels (cholinergic). Sympathetic effects: pupil dilation (mydriasis), increased heart rate and contractility (chronotropic and inotropic—\(\beta\)1), bronchodilation (\(\beta\)2), vasoconstriction of skin and splanchnic vessels (\(\alpha\)1), vasodilation of skeletal muscle vessels (\(\beta\)2), increased sweating (cholinergic), increased renin secretion (\(\beta\)1), decreased GI motility and secretion (\(\alpha\)2, \(\beta\)2), contraction of sphincters (urinary, anal—\(\alpha\)1), ejaculation, and release of glucose (glycogenolysis—\(\beta\)2, gluconeogenesis). The adrenal medulla (modified sympathetic ganglion) secretes adrenaline (epinephrine, 80%) and noradrenaline (20%) into the bloodstream. The sympathetic nervous system also controls blood pressure via the baroreceptor reflex.

\medterm{Synapse}-- The specialised junction between a neuron and another neuron (or effector cell—muscle, gland) where information is transmitted from one cell to another. Types: (1) Chemical synapses (most common in the nervous system)—the presynaptic terminal releases neurotransmitter molecules into the synaptic cleft (20--40 nm wide) in response to an action potential. Neurotransmitters diffuse across the cleft and bind to postsynaptic receptors (ionotropic—ligand-gated ion channels, fast response; metabotropic—G-protein coupled receptors, slow, modulatory). This generates either an excitatory postsynaptic potential (EPSP—depolarisation) or inhibitory postsynaptic potential (IPSP—hyperpolarisation). Summation of many EPSPs reaching threshold generates a new action potential. (2) Electrical synapses—gap junctions (connexin proteins) allow direct flow of ions between cells, enabling very fast, bidirectional transmission (found in smooth and cardiac muscle, some CNS regions). Key neurotransmitters: glutamate (main excitatory CNS transmitter), GABA (main inhibitory CNS transmitter), acetylcholine (neuromuscular junction, autonomic ganglia), dopamine, noradrenaline, serotonin, and neuropeptides (substance P, endorphins). Synaptic plasticity: long-term potentiation (LTP) and long-term depression (LTD)—mechanisms of learning and memory. The synapse is the site of action of many drugs (anaesthetics, antidepressants, antipsychotics, anxiolytics, muscle relaxants) and toxins (botulinum toxin—blocks ACh release, tetanus toxin—blocks inhibitory interneurons).

\medterm{Taste}-- The sensation produced by the interaction of chemical substances with taste receptor cells (TRCs) on the tongue, palate, pharynx, epiglottis, and upper oesophagus. Taste is one of the five primary senses (gustation). Five basic taste qualities: sweet (sugars, artificial sweeteners—T1R2/T1R3 receptors), sour (acids—protons, PKD2L1 receptor), salty (sodium ions—ENaC channel), bitter (alkaloids, toxins—T2R receptor family, ~25 subtypes), and umami (savoury—glutamate, T1R1/T1R3 receptor). Taste buds (~5,000) are housed in papillae on the tongue: fungiform (anterior), foliate (lateral), circumvallate (posterior), and filiform (no taste buds—mechanical). Each taste bud contains 50--100 TRCs with microvilli exposed to the oral cavity via a taste pore. Taste information is transmitted via cranial nerves: facial (CN VII—chorda tympani, anterior two-thirds of tongue), glossopharyngeal (CN IX—posterior one-third), and vagus (CN X—epiglottis and pharynx) to the nucleus of the solitary tract in the medulla, then to the thalamus (VPM) and primary gustatory cortex (insula and operculum). Taste disorders: ageusia (loss of taste), hypogeusia (reduced), dysgeusia (distorted), and phantom taste perception. Causes: viral infections (COVID-19—anosmia and ageusia are hallmark symptoms), medications (ACE inhibitors—metallic taste), chemotherapy, head trauma, radiation, Bell palsy, and ageing. Taste is closely linked to smell (olfaction); much of what is perceived as 'taste' is actually retronasal olfaction.

\medterm{Tay-Sachs Disease}-- A lysosomal storage disease caused by deficiency of hexosaminidase A (Hex A), an autosomal recessive disorder leading to accumulation of GM2 ganglioside in neurons of the central nervous system. Hex A is a heterodimer of alpha and beta subunits; mutations in the HEXA gene (encoding the alpha subunit) cause the classic infantile form. Normal carrier screening: Hex A activity is approximately 50 percent of normal. Affected infants appear normal at birth but develop progressive neurodegeneration by 3-6 months with hypotonia, loss of motor skills, exaggerated startle reflex, cherry-red macula on fundoscopy, macrocephaly, and seizures, with death typically by age 4. Carrier frequency is approximately 1 in 30 in Ashkenazi Jewish populations, and enzyme-based carrier screening has reduced the incidence through prenatal diagnosis.

\medterm{TCA Cycle}-- The tricarboxylic acid (Krebs) cycle occurs in the mitochondrial matrix, completing the oxidation of acetyl-CoA (from carbohydrates, fats, and proteins) to CO2 while generating reducing equivalents (NADH, FADH2) for the electron transport chain. Eight steps: (1) citrate synthase (acetyl-CoA + oxaloacetate to citrate), (2) aconitase (citrate to isocitrate), (3) isocitrate dehydrogenase (rate-limiting, isocitrate to alpha-ketoglutarate, producing NADH and CO2), (4) alpha-ketoglutarate dehydrogenase complex (alpha-ketoglutarate to succinyl-CoA, producing NADH and CO2, requiring TPP, lipoic acid, CoA, FAD, NAD+), (5) succinyl-CoA synthetase (succinyl-CoA to succinate, producing GTP), (6) succinate dehydrogenase (succinate to fumarate, producing FADH2, also Complex II of ETC), (7) fumarase (fumarate to malate), (8) malate dehydrogenase (malate to oxaloacetate, producing NADH). The cycle produces 3 NADH, 1 FADH2, and 1 GTP per acetyl-CoA, yielding approximately 10 ATP equivalents.

\medterm{Transcription}-- The synthesis of RNA from a DNA template by RNA polymerase, the first step of gene expression. In eukaryotes, RNA polymerase II transcribes protein-coding genes into pre-mRNA. Initiation: RNA polymerase binds to the promoter region (containing the TATA box approximately 25-30 base pairs upstream of the transcription start site), with the help of general transcription factors (TFIIA, TFIIB, TFIID/TBP, TFIIE, TFIIF, TFIIH) forming the pre-initiation complex. TFIID contains the TATA-binding protein (TBP) that binds the TATA box. Elongation: RNA polymerase II moves along the template strand, synthesising pre-mRNA in the 5' to 3' direction, adding ribonucleotides complementary to the DNA template. Termination: polyadenylation signal (AAUAAA) triggers cleavage and polyadenylation of the 3' end, followed by RNA polymerase dissociation. Transcription is regulated by promoter elements, enhancers, silencers, and transcription factors, allowing precise spatiotemporal gene expression.

\medterm{Translation}-- The synthesis of proteins from mRNA templates by ribosomes, the final step of gene expression. The genetic code is read in triplets (codons), each specifying an amino acid. Initiation: the small ribosomal subunit (40S in eukaryotes) binds the mRNA at the 5' cap, scans for the start codon (AUG), and the initiator tRNA (Met-tRNAi) binds. The large subunit (60S) joins to form the complete ribosome (80S), with the start codon in the P site. Eukaryotic initiation factors (eIFs) facilitate this process: eIF4E binds the 5' cap, eIF2 binds Met-tRNAi to the small subunit, and eIF5B joins the large subunit. Elongation: aminoacyl-tRNAs enter the A site via EF-Tu (eEF1A in eukaryotes), peptide bond formation occurs in the P site, and translocation moves the ribosome three nucleotides along the mRNA via EF-G (eEF2). Termination: a stop codon (UAA, UAG, UGA) in the A site is recognised by release factors, releasing the completed polypeptide.

\medterm{Tubuloglomerular Feedback}-- A feedback mechanism at the juxtaglomerular apparatus (JGA) that adjusts GFR based on tubular fluid NaCl concentration. The macula densa cells (specialized tubular epithelial cells in the thick ascending limb) sense NaCl concentration via the NKCC2 cotransporter. When GFR is high, increased NaCl delivery to the macula densa triggers release of ATP and adenosine, causing afferent arteriole constriction and reducing GFR. When GFR is low, decreased NaCl delivery causes afferent arteriole dilation (via increased nitric oxide and prostaglandins), increasing GFR. This feedback loop maintains a relatively constant NaCl delivery to the distal tubule and stabilises GFR despite changes in blood pressure. The JGA also releases renin in response to low NaCl, activating the RAAS. Tubuloglomerular feedback is impaired in diabetic nephropathy and during ACE inhibitor therapy.

\medterm{Uncompetitive Inhibition}-- A type of enzyme inhibition where the inhibitor binds only to the enzyme-substrate (ES) complex at a site distinct from the active site, forming an inactive ESI complex. The inhibitor cannot bind to the free enzyme. Both Vmax and Km decrease: Vmax decreases because the ESI complex is non-productive (fewer active ES complexes), and Km decreases because the inhibitor effectively traps substrate in the ES complex (increasing apparent substrate affinity). On a Lineweaver-Burk plot: lines are parallel, with different slopes and both x- and y-intercepts changing (same slope, different y-intercept and x-intercept). Uncompetitive inhibition is rare in single-substrate reactions but occurs in multisubstrate reactions (e.g., inhibition of lactate dehydrogenase by pyruvate at high concentrations) and is the mechanism of some drugs such as methotrexate (inhibition of dihydrofolate reductase when bound to NADPH).

\medterm{Urea Cycle}-- The hepatic metabolic pathway converting toxic ammonia (NH3/NH4+) to urea for renal excretion, occurring partly in the mitochondrial matrix and partly in the cytoplasm. Five steps: (1) carbamoyl phosphate synthetase I (CPS-I, mitochondrial, rate-limiting, activated by N-acetylglutamate), combining CO2, NH3, and ATP; (2) ornithine transcarbamylase (OTC, mitochondrial), combining carbamoyl phosphate with ornithine to form citrulline; (3) argininosuccinate synthetase (cytoplasmic), combining citrulline with aspartate to form argininosuccinate; (4) argininosuccinate lyase (cytoplasmic), cleaving argininosuccinate to arginine and fumarate; (5) arginase (cytoplasmic), hydrolysing arginine to ornithine and urea. Urea is the non-toxic nitrogenous waste product excreted by the kidneys. Deficiencies cause hyperammonaemia (ammonia intoxication), with symptoms including lethargy, vomiting, cerebral oedema, and coma.

\medterm{Vitamin A: Retinal}-- Vitamin A (retinol) is a fat-soluble vitamin essential for vision, immune function, cell differentiation, and reproduction. In the eye, retinol is oxidized to retinal (retinaldehyde), which binds to opsin in the photoreceptor outer segments forming rhodopsin (in rods) and photopsins (in cones). Light absorption by 11-cis-retinal causes isomerization to all-trans-retinal, triggering the phototransduction cascade (G-protein transducin activation, phosphodiesterase activation, cGMP hydrolysis, Na+ channel closure, and hyperpolarisation of the photoreceptor cell). Vitamin A deficiency causes night blindness (nyctalopia), xerophthalmia (dry cornea), Bitot spots, keratomalacia, and increased infection risk. Excess vitamin A (hypervitaminosis A) causes hepatotoxicity, teratogenicity, and raised intracranial pressure.

\medterm{Vitamin B Complex}-- The vitamin B complex comprises eight water-soluble B vitamins that function primarily as coenzymes or cofactors in cellular metabolism: Vitamin B1 (thiamine, a cofactor in carbohydrate metabolism and nerve function), Vitamin B2 (riboflavin, precursor of FAD/FMN in oxidation-reduction reactions), Vitamin B3 (niacin, precursor of NAD+/NADP+), Vitamin B5 (pantothenic acid, component of coenzyme A), Vitamin B6 (pyridoxine, cofactor in amino acid metabolism and neurotransmitter synthesis), Vitamin B7 (biotin, cofactor in carboxylation reactions), Vitamin B9 (folate/folic acid, essential for one-carbon transfer in nucleotide synthesis), and Vitamin B12 (cobalamin, cofactor for methionine synthase and methylmalonyl-CoA mutase). B vitamins are obtained from a wide variety of foods (meat, eggs, dairy, legumes, leafy greens, whole grains) and are generally not stored in large quantities, requiring regular dietary intake. Deficiencies cause specific syndromes (beriberi for B1, pellagra for B3, megaloblastic anaemia for B9/B12, see individual entries).

\medterm{Vitamin B12: Cobalamin}-- Vitamin B12 (cobalamin) is a water-soluble vitamin containing cobalt, essential for two enzymatic reactions: (1) methionine synthase (cytoplasmic, requiring methylcobalamin), converting homocysteine to methionine (linking to folate metabolism), and (2) methylmalonyl-CoA mutase (mitochondrial, requiring adenosylcobalamin), converting methylmalonyl-CoA to succinyl-CoA (in odd-chain fatty acid and branched-chain amino acid metabolism). Intrinsic factor (IF, from gastric parietal cells) is essential for B12 absorption in the terminal ileum. B12 deficiency causes megaloblastic anaemia (with hypersegmented neutrophils), peripheral neuropathy, subacute combined degeneration of the spinal cord (dorsal columns and corticospinal tracts), and glossitis. Pernicious anaemia is an autoimmune condition causing IF deficiency.

\medterm{Vitamin B1: Thiamine}-- Thiamine (thiamine pyrophosphate, TPP) is a water-soluble B vitamin essential as a cofactor for several key metabolic enzymes: pyruvate dehydrogenase (linking glycolysis to the TCA cycle), alpha-ketoglutarate dehydrogenase (TCA cycle), transketolase (pentose phosphate pathway), and branched-chain alpha-ketoacid dehydrogenase (branched-chain amino acid catabolism). Thiamine deficiency causes: beriberi (dry: peripheral neuropathy and muscle wasting; wet: high-output heart failure with edema and vasodilation), and Wernicke-Korsakoff syndrome (confusion, ataxia, ophthalmoplegia, and amnesia, commonly seen in chronic alcoholism). Wernicke encephalopathy is a medical emergency treated with intravenous thiamine before glucose administration to prevent Korsakoff psychosis.

\medterm{Vitamin B3: Niacin}-- Niacin (nicotinic acid and nicotinamide) is a water-soluble vitamin that serves as the precursor for the coenzymes NAD+ (nicotinamide adenine dinucleotide) and NADP+ (nicotinamide adenine dinucleotide phosphate), which are essential for over 400 enzymatic reactions including glycolysis, the TCA cycle, oxidative phosphorylation, and fatty acid synthesis. NAD+ is a key electron carrier in catabolic reactions; NADP+ is important in anabolic reactions and antioxidant defense (via glutathione). Niacin deficiency causes pellagra, characterised by the 4 Ds: dermatitis (photosensitive), diarrhoea, dementia, and death. Pellagra is endemic in populations consuming maize-based diets low in tryptophan (niacin precursor). Therapeutic niacin (in high doses) lowers triglycerides and LDL while raising HDL.

\medterm{Vitamin B6: Pyridoxine}-- Pyridoxine (pyridoxal phosphate, PLP, the active form) is a water-soluble B vitamin essential as a cofactor for over 100 enzymatic reactions, particularly in amino acid metabolism. Key reactions include transamination (ALT, AST), decarboxylation (producing neurotransmitters: serotonin from 5-HTP, dopamine/norepinephrine from L-DOPA, GABA from glutamate, histamine from histidine), and glycogen phosphorylase (glycogenolysis). PLP is also required for heme synthesis (ALA synthase step), sphingolipid synthesis, and serine racemization. PLP is also required for amino acid transport across cell membranes. Pyridoxine deficiency is rare but causes dermatitis, glossitis, microcytic anaemia, peripheral neuropathy, and seizures in infants. Isoniazid (tuberculosis treatment) can cause B6 deficiency as a side effect, requiring supplementation.

\medterm{Vitamin B9: Folate}-- Folate (folic acid, tetrahydrofolate/THF) is a water-soluble B vitamin essential for one-carbon transfer reactions in nucleotide synthesis (thymidylate synthesis for DNA, purine synthesis), amino acid metabolism (methionine synthesis from homocysteine), and methylation reactions. THF carries one-carbon units at various oxidation states (methyl, methylene, methenyl, formyl, formimino). Key reactions: thymidylate synthase converts dUMP to dTMP (DNA synthesis, the target of 5-fluorouracil), and methionine synthase converts homocysteine to methionine (requiring methylcobalamin/B12). Folate deficiency causes megaloblastic anaemia (identical to B12 deficiency), neural tube defects in pregnancy (spina bifida, anencephaly), and elevated homocysteine (cardiovascular risk factor). Folic acid supplementation before and during early pregnancy prevents neural tube defects.

\medterm{Vitamin C (Ascorbic Acid)}-- Vitamin C (ascorbic acid) is a water-soluble vitamin and a powerful antioxidant, essential for collagen synthesis (required for proline and lysine hydroxylation in collagen triple helix formation), carnitine synthesis, neurotransmitter synthesis (dopamine to norepinephrine), and iron absorption (reduces ferric to ferrous iron in the gut). Humans lack L-gulonolactone oxidase and cannot synthesise vitamin C, requiring dietary intake (citrus fruits, berries, kiwi, tomatoes, peppers, broccoli). Vitamin C deficiency causes scurvy: impaired wound healing, gingival hyperplasia and bleeding (swollen, bleeding gums), perifollicular haemorrhages, ecchymoses, coiled hairs, joint pain, fatigue, and anaemia. In severe cases, scurvy causes haemarthrosis, intraosseous bleeding, and death. Vitamin C is also a cofactor for dopamine beta-hydroxylase (norepinephrine synthesis). High-dose vitamin C has been used as an unproven treatment for the common cold and has been investigated in sepsis and cancer (intravenous high-dose as adjunctive therapy, limited evidence).

\medterm{Vitamin D: Calcitriol}-- Vitamin D is a fat-soluble vitamin (actually a prohormone) essential for calcium and phosphorus homeostasis and bone health. Cholecalciferol (D3) is synthesized in the skin from 7-dehydrocholesterol by UVB radiation, or obtained from diet (ergocalciferol/D2 from plants, D3 from animal sources). It is hydroxylated in the liver to 25-hydroxyvitamin D (calcidiol), then in the kidney to the active form 1,25-dihydroxyvitamin D (calcitriol) by 1-alpha-hydroxylase (regulated by PTH, low calcium, and low phosphate; inhibited by FGF23 and high phosphate). Calcitriol increases intestinal calcium and phosphate absorption, promotes renal calcium reabsorption, and mobilises calcium from bone. Vitamin D deficiency causes rickets in children (soft, deformed bones) and osteomalacia in adults (bone pain, weakness, fractures).

\medterm{Vitamin E: Tocopherol}-- Vitamin E (alpha-tocopherol) is a fat-soluble vitamin and the body's primary lipid-soluble antioxidant, protecting polyunsaturated fatty acids in cell membranes from lipid peroxidation by free radicals (reactive oxygen species). It donates a hydrogen atom to lipid peroxyl radicals, breaking the chain reaction of lipid peroxidation. Vitamin E works synergistically with vitamin C (which regenerates reduced vitamin E at the membrane-cytosol interface) and selenium (as a cofactor for glutathione peroxidase) in protecting against oxidative damage. Vitamin E deficiency is rare but occurs in fat malabsorption syndromes (cystic fibrosis, abetalipoproteinaemia) and causes peripheral neuropathy, spinocerebellar ataxia, retinopathy, and haemolysis (due to red blood cell membrane lipid peroxidation).

\medterm{Vitamin K: Phylloquinone}-- Vitamin K is a fat-soluble vitamin essential for the post-translational gamma-carboxylation of glutamate residues to gamma-carboxyglutamate (Gla) in specific clotting factors (II, VII, IX, X, protein C, protein S) and anticoagulant proteins (protein C, protein S, protein Z). The gamma-carboxylation reaction, catalyzed by gamma-glutamyl carboxylase, requires reduced vitamin K (vitamin K hydroquinone) as a cofactor and is essential for the calcium-binding function of these proteins, allowing them to bind calcium and phospholipid membranes. Vitamin K deficiency causes bleeding (elevated PT/INR), most commonly in newborns (given prophylactic vitamin K at birth), patients with fat malabsorption, or those on warfarin therapy (which inhibits vitamin K epoxide reductase, preventing recycling of vitamin K). The antidote is vitamin K administration.

\medterm{Zymogen (Proenzyme)}-- An inactive precursor of an enzyme that requires proteolytic cleavage to become catalytically active. Zymogens are synthesised and stored in cells (typically in secretory vesicles or zymogen granules) and are activated at the appropriate site and time by specific proteases, preventing premature or inappropriate enzymatic activity that could damage the cell of origin or surrounding tissues. The conversion of a zymogen to an active enzyme involves the cleavage of one or more specific peptide bonds, which induces a conformational change that exposes the active site. Major zymogen systems in human physiology: digestive zymogens (secreted by the exocrine pancreas and gastric chief cells as inactive precursors to prevent autodigestion — trypsinogen → trypsin, chymotrypsinogen → chymotrypsin, proelastase → elastase, procarboxypeptidase → carboxypeptidase, pepsinogen → pepsin in the stomach, activated by gastric acid), clotting cascade (coagulation factors circulate as inactive zymogens — prothrombin → thrombin, fibrinogen → fibrin, factor X → Xa, etc. — activated sequentially in the coagulation cascade), fibrinolytic system (plasminogen → plasmin), complement system (complement proteins C1-C9 circulate as inactive zymogens, activated by classical, lectin, or alternative pathways), and procollagen (procollagen → collagen by procollagen peptidases). The pancreas is particularly dependent on zymogen physiology: the pancreatic acinar cells synthesise and package digestive enzymes as zymogens in zymogen granules, which are secreted into the pancreatic duct and activated in the duodenum by enteropeptidase (enterokinase). Premature activation of trypsinogen to trypsin within the pancreas causes acute pancreatitis (autodigestion of pancreatic tissue)., essential for collagen synthesis (required for proline and lysine hydroxylation in collagen triple helix formation), carnitine synthesis, neurotransmitter synthesis (dopamine to norepinephrine), and iron absorption (reduces ferric to ferrous iron in the gut). Humans lack L-gulonolactone oxidase and cannot synthesise vitamin C, requiring dietary intake (citrus fruits, berries, kiwi, tomatoes, peppers, broccoli). Vitamin C deficiency causes scurvy: impaired wound healing, gingival hyperplasia and bleeding (swollen, bleeding gums), perifollicular haemorrhages, ecchymoses, coiled hairs, joint pain, fatigue, and anaemia. In severe cases, scurvy causes haemarthrosis, intraosseous bleeding, and death. Vitamin C is also a cofactor for dopamine beta-hydroxylase (norepinephrine synthesis). High-dose vitamin C has been used as an unproven treatment for the common cold and has been investigated in sepsis and cancer (intravenous high-dose as adjunctive therapy, limited evidence).
